\itemhead{习题4.1 · 第13题（教材p.105）}{XJ-B1-C04-HW-P105-Y0620}
\label{XJ-B1-C04-HW-P105-Y0620}
{\footnotesize 湘教版教材p.105 / PDF第112页 / 习题4.1 / 第13题}
\field{原题}{已知幂函数 $f(x)=x^{2m^2-m-6}\ (m\in\mathbb Z)$ 在 $(0,+\infty)$ 上是减函数。
\subq{1}求 $f(x)$ 的表达式；
\subq{2}判断 $f(x)$ 的奇偶性和单调性；
\subq{3}求 $f(x)$ 的值域。}
\field{参考答案}{$2m^2-m-6<0$ 得 $-3/2<m<2$，整数 $m=-1,0,1$ 分别给指数 $-3,-6,-5$。$x^{-3},x^{-5}$ 为奇函数，在两半轴分别递减，值域 $\mathbb R\setminus\{0\}$；$x^{-6}$ 为偶函数，在负半轴递增、正半轴递减，值域 $(0,+\infty)$。三者定义域均排除0。}
{\small \field{关联知识}{P03 幂函数的概念、图象与性质（教材帮p.125、126、127 / PDF134、135、136）；F08 奇偶性、对称与函数性质应用（教材帮p.113、114、115 / PDF122、123、124）；F05 单调性定义、证明与参数（教材帮p.97、98、100 / PDF106、107、109）}
\field{方法与易错点}{设 $f(x)=x^\alpha$，用已知点求指数，再查该指数的定义域、奇偶性和各单调区间。 易错：把 $3x^2$ 当标准幂函数；将第一象限性质直接用于负半轴；分数负指数遗漏0禁点。}
\field{答案依据}{答案册相关题组：PDF第31页（册内页30）；按当前题库原题逐项重算，不按旧版题号机械对套。}}
